\documentclass[11pt]{article}
\usepackage{icmlko}
\begin{document}

\title{잘못된 자리를 가리킨 후보였는데 재 보니 답이 나왔다 --- 위키도 닫는다}
\author{월드모델 실험실}
\date{노트 596 · 2026년 8월 5일}
\maketitle

\begin{abstract}
노트 $548$ 이 \textbf{검색 재수집을 닫았다} --- 축 전체를 가려 판이 얼마나
내려가나로 \textbf{상금 상한}을 재고(판 $\mathbf{0.0204}$), 재수집의
기대값이 그보다 훨씬 작아 안 하기로 했다. \textbf{위키는 같은 자로 안 재
봤다.}

\begin{center}
\begin{tabular}{lrl}
\hline
 & 상금 상한 & \\
\hline
검색(노트 $548$) & $0.0204$ & 닫음 \\
\textbf{위키(여기)} & $\mathbf{0.0086}$ & \textbf{검색의 $42\%$ --- 닫음} \\
\hline
\end{tabular}
\end{center}

그런데 이 사이클의 값은 그 숫자가 아니라 \textbf{내가 세운 전제 둘이 깨진
것}이다.
\end{abstract}

\begin{figure}[t]
\centering
\includegraphics[width=\textwidth]{figs/f1.pdf}
\caption{\textbf{가} --- 크게 잃는 곳이 작은 도메인이다.
\textbf{나} --- 검색의 절반도 안 된다.
\textbf{다} --- 내 전제 둘이 깨졌다.}
\end{figure}

\section{전제 ① --- 위키가 어디 있는지 틀렸다}

후보를 세울 때 노트 $554$ 의 ``위키도 같다(만화 $16.6\%$ $\cdot$ 펀딩
$13.8\%$ 만 조회)''를 \textbf{하네스 덮음}으로 읽었다. 실제로 재니

\begin{center}
\begin{tabular}{lrrr}
\hline
도메인 & 덮음 & 유보 & 위키를 빼면 \\
\hline
\textbf{아이돌} & $46\%$ & $51$ & $\mathbf{-0.0730}$ \\
\textbf{세계애니} & $42\%$ & $300$ & $\mathbf{-0.0611}$ \\
게임 & $29\%$ & $180$ & $-0.0339$ \\
모바일 & $6\%$ & $441$ & $+0.0002$ \\
만화 & $\mathbf{0.5\%}$ & $258$ & $+0.0016$ \\
펀딩 & $\mathbf{0\%}$ & $529$ & --- \\
\hline
\end{tabular}
\end{center}

\textbf{만화 $0.5\%$ $\cdot$ 펀딩 $0\%$.} 노트 $554$ 의 수치는
\emph{원천 조회율}이었고 하네스 덮음이 아니었다. \textbf{내 후보는 위키가
없는 두 도메인을 가리키고 있었다.}

\begin{quote}
\textbf{그래도 재 보니 답이 나왔다.} 후보가 틀린 자리를 가리켰어도
\emph{그 축 전체}를 가리는 팔은 어차피 옳은 물음이었다 ---
``지금 있는 위키가 판에 얼마를 주나''.
\end{quote}

\section{전제 ② --- 집 밖이 움직였다}

배선 가드로 ``집 밖 짝은 위키를 안 받으므로(ALL5 만) 안 움직여야 한다''를
걸었다. \textbf{KR 만화 $-0.0045$($0$ 밖 음수 $\cdot$ $4/12$)} 다.

\begin{quote}
\textbf{이웃 효과다}(노트 $551$) --- 위키를 빼면 \emph{공유 모형}이 바뀌고
그것이 집 밖 예보에 실린다. 앱은 $+0.0013$ 으로 $0$ 을 문다. 크기가 작고
방향이 갈려 결론은 안 바뀌지만, \textbf{``짝이 그 축을 안 받으니 안
움직인다''는 틀린 전제였다.}
\end{quote}

이 판이 이웃 효과를 여섯 번째로 만난다(노트 $525 \cdot 551 \cdot 559
\cdot 560 \cdot 563 \cdot$ 여기). \textbf{매번 다른 옷을 입고 온다.}

\section{왜 상한이 작나}

도메인별로는 크게 잃는다 --- 아이돌 $-0.073$, 세계애니 $-0.061$, 게임
$-0.034$. 그런데 \textbf{위키가 사는 곳이 작은 도메인이다.}

\begin{center}
$51 + 300 + 180 = 531$ 행 $=$ 분모 $3{,}369$ 의 $\mathbf{16\%}$
\end{center}

\textbf{도메인별 $-0.07$ 이 판에서 $-0.009$ 가 된다.} 노트 $548$ 이 검색을
닫을 때와 같은 셈이고, 여기서는 더 확실하다.

\section*{조항}

\begin{quote}
\textbf{후보를 세운 전제가 틀려도 실험은 옳을 수 있다.} 잘못된 자리를
가리킨 후보였지만, 팔이 ``\emph{그 축 전체}를 가린다''였으므로 물음 자체는
맞았다. \textbf{후보의 근거와 팔의 설계는 따로 검사해야 한다} ---
근거가 틀렸다고 결과를 버리면 안 되고, 근거가 맞았다고 팔을 안 봐도 되는
것이 아니다.
\end{quote}

\section*{남은 것}

\begin{itemize}
\item \textbf{아이돌 $\cdot$ 세계애니는 위키에 크게 기댄다}($-0.073 \cdot
-0.061$). 그 둘의 덮음을 $46\% \cdot 42\%$ 에서 올리면 그 도메인은 오를
수 있다 --- 다만 무게가 작아(합쳐 $351$ 행, $10\%$) \textbf{판 이득의
상한이 $0.006$ 쯤}이다. 노트 $548$ 규칙대로 \textbf{안 쫓는다.}
\item \textbf{펀딩이 왜 $0\%$ 인가} --- \texttt{WIKI\_DROP} 에 없는데 덮음이
$0$ 이다. 수집을 안 했거나 붙지 않았다. 안 봤다.
\item \textbf{이웃 효과의 크기를 예측할 수 없다} --- 여섯 번을 만났는데
매번 사후에만 안다. 그것을 미리 재는 자가 없다.
\end{itemize}

\end{document}
